% Generated draft: requires mathematical and source review.
\documentclass[11pt]{article}
\usepackage[margin=25mm]{geometry}
\usepackage{fontspec}
\setmainfont{Latin Modern Roman}
\usepackage{amsmath,amssymb,mathtools,unicode-math}
\setmathfont{Latin Modern Math}
\usepackage{enumitem,needspace,xurl,hyperref}
\setlength{\parindent}{0pt}
\setlength{\parskip}{5pt}
\setlength{\emergencystretch}{3em}
\newcommand{\N}{\mathbb{N}}
\newcommand{\Z}{\mathbb{Z}}
\newcommand{\Q}{\mathbb{Q}}
\newcommand{\R}{\mathbb{R}}
\newcommand{\C}{\mathbb{C}}
\newcommand{\F}{\mathbb{F}}
\newcommand{\A}{\mathbb{A}}
\newcommand{\PP}{\mathbb{P}}
\newcommand{\E}{\mathbb{E}}
\newcommand{\eps}{\varepsilon}
\newcommand{\dd}{\,\mathrm{d}}
\begin{document}
\section*{UnsolvedMath ID: 30003361}
\textbf{Source problem code:} catalogue entry
\subsection{Definitions and mathematical statement}
For context, if \(M\) is a smooth manifold, let \(\operatorname{Riem}(M)\) denote the set of smooth positive-definite symmetric covariant \(2\)-tensor fields on \(M\), and let \(\operatorname{Diff}(M)\) denote its group of smooth diffeomorphisms. The usual equivalence relation on metrics is
\[
 g_1\sim g_2\quad\Longleftrightarrow\quad g_1=\phi^*g_2\text{ for some }\phi\in\operatorname{Diff}(M),
\]
with quotient \(\operatorname{Riem}(M)/\operatorname{Diff}(M)\). This is contextual notation; the source does not specify a manifold, dimension, compactness assumption, or boundary condition.

The supplied research question is: \emph{Can diffeomorphisms be replaced by homeomorphisms and Riemannian metrics by singular metrics in the study of moduli spaces?}

A schematic formulation would require a specified underlying topological space \(X\), a specified class \(\mathcal S(X)\) of singular metric structures, and a rule for pulling these structures back by every element of \(\operatorname{Homeo}(X)\), the group of self-homeomorphisms of \(X\). If such a rule is defined and preserves \(\mathcal S(X)\), the proposed equivalence relation would be
\[
 s\sim t\quad\Longleftrightarrow\quad s=h^*t\text{ for some }h\in\operatorname{Homeo}(X),
 \qquad s,t\in\mathcal S(X),
\]
and the corresponding quotient would be \(\mathcal S(X)/\operatorname{Homeo}(X)\).

The question asks whether an appropriate moduli theory can be developed in this setting. The source does not define singular metrics, the admissible spaces \(X\), the pullback rule, the relevant topologies, or the properties that such a theory must establish. In particular, existence of the displayed orbit set alone does not specify the intended mathematical goal. These choices remain unresolved, so the supplied question is a research programme rather than a fully quantified conjecture.
\subsection{Short English statement}
Can the study of metric moduli spaces be extended to singular metric structures, with equivalence given by homeomorphisms? The source leaves the allowed singularities, underlying spaces, topology, and desired conclusions unspecified, preventing a unique precise conjecture.
\subsection{Sources}
\begin{itemize}
\item Supplied local UnsolvedMath registry record, permanent ID \(30003361\), title \emph{Moduli of Singular Metrics Under Homeomorphisms}, source problem code \texttt{catalogue entry}, in the collection \emph{Oberwolfach Reports Open Problems}. The registry records proposed year 2017 and names no individual proposer. Its statement status is supplied as exact, with the qualification that the wording was retained byte-for-byte after an internal readability check; this is not independent source verification.
\item Supplied report citation: \emph{Mini-Workshop: Spaces and Moduli Spaces of Riemannian Metrics (2018)}, Oberwolfach Reports, pp.\ 133--166, DOI \url{https://doi.org/10.4171/owr/2017/3}. The background attributes the problem to workshop open-problem discussions documented in Oberwolfach Reports. The citation's 2018 designation and the record's proposed year 2017 are retained without resolving their relationship.
\item Supplied literature source: \emph{Oberwolfach discussion of singular metric moduli}, \url{https://doi.org/10.4171/owr/2017/3}.
\item Supplied literature source: \emph{Positive scalar curvature on manifolds with singularities}, \url{https://doi.org/10.3842/SIGMA.2021.034}.
\item Supplied literature source: \emph{Moduli spaces of Riemannian metrics}, \url{https://doi.org/10.1007/978-3-0348-0948-1}.
\item Literature assessment dated \texttt{2026-08-21}: the supplied curation labels the problem open and reports partial theories for observer quotients, Alexandrov or CAT-type spaces, and singular positive-scalar-curvature constructions, while reporting no general theory making the proposed replacements. This is supplied dated curation derived from the OWR material, not an independently verified literature assessment. The source explicitly requests independent verification of time-sensitive status and cited claims.
\end{itemize}
\end{document}
